% --- Back cover ---
% Large PERSONIX wordmark, no logo image. Cream paper inherited globally.
% Absolute (tikz overlay) positioning so the footer sticks to the bottom.
\clearpage
\thispagestyle{empty}
\begin{tikzpicture}[remember picture, overlay]
  % --- Wordmark ---
  \node[anchor=north, font=\sffamily\bfseries\fontsize{48}{52}\selectfont]
    at ([yshift=-26mm]current page.north) {PERSONIX};
  \node[anchor=north, font=\rmfamily\itshape\large]
    at ([yshift=-44mm]current page.north) {تغيير بلا تنازلات};
  \node[anchor=north, font=\rmfamily\small,
        text width=\paperwidth-50mm, align=center]
    at ([yshift=-54mm]current page.north)
    {شبكة سمعة لامركزية لا يمكن الرقابة عليها ولا إفسادها};
  % --- Body ---
  \node[anchor=north, text width=\paperwidth-50mm, align=justify, font=\small]
    at ([yshift=-72mm]current page.north)
    {%
      نجحت الدولة حين كان التواصل بطيئًا، والقرارات محلية، والسلطة تعيش في المدينة
      نفسها التي تحكمها. أما شبكات اليوم فتقلب هذه الفرضيات الثلاث جميعًا --- ولم تعد
      المؤسسات المبنية عليها تلائم العالم الذي تحكمه. يصف هذا الكتاب أداة تلائمه:
      شبكة سمعة لامركزية تبقى فيها الهوية مِلكك، والحقيقة قابلة للتدقيق علنًا،
      والرشوة انتحار سمعي.

      \smallskip
      ليس بيانًا حماسيًا. ولا تنبؤًا. بل مخططًا عمليًا لكيفية بناء خليفة الدولة ---
      طوعًا، إعلانًا تلو إعلان.
    };
  % --- Footer (anchored to the bottom of the page) ---
  \node[anchor=south, font=\sffamily\footnotesize,
        text width=\paperwidth-50mm, align=center]
    at ([yshift=12mm]current page.south)
    {%
      Pavel Kudrna\quad\textbullet\quad براغ، 2026\par
      \href{https://personix.org}{personix.org}\par
      مرخّص بموجب Creative Commons Attribution-ShareAlike 4.0 International
      (CC BY-SA 4.0).
    };
\end{tikzpicture}%
\mbox{}%
\clearpage
